\chapter{Boolean Analysis — Cutoff}
%%% generated by the blueprint pipeline from TCSlib.BooleanAnalysis.GaussianSpace.Invariance.Cutoff
\section{Overview}

Declarations of the Lean module \texttt{TCSlib.BooleanAnalysis.GaussianSpace.Invariance.Cutoff}.

\section{Declarations}

\begin{definition}[Cutoff ramp]
\lean{BooleanAnalysis.GaussianSpace.Invariance.cutoffRamp}
\label{BooleanAnalysis.GaussianSpace.Invariance.cutoffRamp}
\leanok
For real numbers $v$, $\eta$, and $x$, define the cutoff ramp by
$r(v,\eta,x)=\max\{0,\min\{1,(v-x)/\eta\}\}$.
\end{definition}

\begin{theorem}[Lipschitz bound for a cutoff ramp]
\lean{BooleanAnalysis.GaussianSpace.Invariance.lipschitz_cutoffRamp}
\label{BooleanAnalysis.GaussianSpace.Invariance.lipschitz_cutoffRamp}
\leanok
\uses{BooleanAnalysis.GaussianSpace.Invariance.cutoffRamp}
\difficulty{3}
For $\beta>0$ and $v\mathbin\bbr$, define $r(x)=\max(0,\min(1,(v-x)/\beta))$ for $x\mathbin\bbr$.
This function is Lipschitz with constant $1/\beta$, meaning $|r(x)-r(y)|\langle |x-y|/\beta$
for all real $x,y$.
\end{theorem}

\begin{theorem}[Bounds for the cutoff ramp]
\lean{BooleanAnalysis.GaussianSpace.Invariance.cutoffRamp_bounds}
\label{BooleanAnalysis.GaussianSpace.Invariance.cutoffRamp_bounds}
\leanok
\uses{BooleanAnalysis.GaussianSpace.Invariance.cutoffRamp}
\difficulty{4}
For real numbers $\eta>0$ and $v$, the cutoff ramp takes values in $[0,1]$ at every real
$x$. It equals $1$ whenever $x+\eta\le v$ and equals $0$ whenever $v\le x$.
\end{theorem}

\begin{theorem}[Smooth cutoff with derivative bounds]
\lean{BooleanAnalysis.GaussianSpace.Invariance.exists_contDiff_cutoff}
\label{BooleanAnalysis.GaussianSpace.Invariance.exists_contDiff_cutoff}
\leanok
\uses{BooleanAnalysis.GaussianSpace.Invariance.cutoffRamp, BooleanAnalysis.GaussianSpace.Invariance.cutoffRamp_bounds, BooleanAnalysis.GaussianSpace.Invariance.intervalSmooth, BooleanAnalysis.GaussianSpace.Invariance.intervalSmooth_iterate_bounds, BooleanAnalysis.GaussianSpace.Invariance.intervalSmooth_iterate_exists_value, BooleanAnalysis.GaussianSpace.Invariance.lipschitz_cutoffRamp}
\difficulty{5}
For every natural number $m$, real number $u$, and real number $\eta>0$, there is a
function $\phi:\mathbb R\to\mathbb R$ of class $C^{m+1}$ taking values in $[0,1]$, equal
to $1$ for $x\le u$ and equal to $0$ for $x\ge u+(m+2)\eta$. Moreover, its $(m+1)$st
derivative satisfies $|\phi^{(m+1)}(x)|\le 2^m/\eta^{m+1}$ for every real $x$.
\end{theorem}
